% BHU PYQ -- Manually transcribed & error-corrected
% Paper: MTB-505 : Mechanics
% Exam: B.A./B.Sc. (Hons) Semester V Examination, 2022-23

\documentclass[11pt,a4paper]{article}
\usepackage[a4paper,top=2.5cm,bottom=2.5cm,left=2.8cm,right=2.8cm,headheight=14pt]{geometry}
\usepackage{amsmath,amssymb}
\usepackage[T1]{fontenc}\usepackage[utf8]{inputenc}
\usepackage{lmodern,microtype,fancyhdr,enumitem,hyperref,lastpage,parskip}

\hypersetup{colorlinks=true,urlcolor=black,linkcolor=black,
  pdftitle={MTB-505: Mechanics -- BHU B.A./B.Sc. Sem V 2022-23}}
\pagestyle{fancy}\fancyhf{}
\renewcommand{\headrulewidth}{0.4pt}\renewcommand{\footrulewidth}{0.4pt}
\fancyhead[L]{\small   \textit{Banaras Hindu University}}
\fancyhead[R]{\small   \textit{MTB-505: Mechanics}}
\fancyfoot[C]{\small Page  \thepage\ of \pageref{LastPage}}


\newlist{parts}{enumerate}{2}
\setlist[parts,1]{label=(\alph*),leftmargin=*,itemsep=5pt,topsep=3pt}
\setlist[parts,2]{label=(\roman*),leftmargin=*,itemsep=3pt,topsep=2pt}

\newcommand{\pts}[1]{\hfill   \textit{[#1]}}


\begin{document}

\begin{center}
  {\large\bfseries BANARAS HINDU UNIVERSITY}\\[2pt]
  {
\normalsize B.A./B.Sc. (Hons) --- Semester V Examination, 2022-23}\\[6pt]
  \rule{0.8\linewidth}{0.5pt}\\[6pt]
  {\large\bfseries Paper No. MTB-505: Mechanics}\\[4pt]
  {
\normalsize Time: 3 Hours \hfill Maximum Marks: 70}
\end{center}
\rule{\linewidth}{0.4pt}
\smallskip



\noindent   \textit{Instructions: Answer   \textbf{five} questions including Question~1 which is compulsory. Figures in right-hand margin indicate marks.}
\bigskip




\noindent   \textbf{Question 1.}

\begin{parts}
  \item Reduce the system of forces $(3,1,-2)$, $(2,-1,-1)$ and $(3,-4,1)$ acting at points $(2,3,-2)$, $(5,0,1)$ and $(-1,3,-1)$, respectively, in the $OXYZ$ frame into a resultant force and a resultant couple. \pts{2}
  \item What are the invariants under the reduction of forces using any base and system of coordinate axes? Give reasons for your answer. \pts{2}
  \item Explain the concept of unstable equilibrium of a heavy body resting on a fixed body. Give an example. \pts{2}
  \item A system consists of two particles of masses $2$ and $3$ units at positions $(1,2,3)$ and $(-1,0,4)$. Find the moment of inertia about the $x$-axis. \pts{2}
  \item The sides of a rectangular solid plate of mass $50\, \text{g}$ are formed by the four lines $y = -3$, $y = 4$, $x = 5$, $x = -4$. Find the product of inertia about the lines $y = 4$ and $x = 5$. \pts{2}
  \item What do you mean by equi-momental systems? \pts{2}
  \item Find the length of the equivalent simple pendulum when the rigid body is a circular disc and the fixed axis is its tangent. \pts{2}
  \item Derive the expression for the kinetic energy of a body rotating about a fixed axis passing through it. \pts{2}
  \item Find the position of the centre of percussion for a uniform rod of length $a$ with one end fixed. \pts{2}
  \item Is the rolling of a circular body on an inclined flat surface under gravity a two-dimensional motion? Give reasons. \pts{2}
  \item State the principle of conservation of linear momentum. \pts{2}
\end{parts}

\medskip\rule{\linewidth}{0.2pt}

\medskip


\noindent   \textbf{Question 2.}

\begin{parts}
  \item Show that any system of forces acting on a rigid body can be reduced to a single force acting at an arbitrarily chosen point and a couple. \pts{6}
  \item Equal forces act along the coordinate axes $OX$, $OY$, $OZ$ and along the line $\dfrac{x}{1} = \dfrac{y}{1} = \dfrac{z}{1}$. Find the equation of the central axis of the system. \pts{6}
\end{parts}

\medskip\rule{\linewidth}{0.2pt}

\medskip


\noindent   \textbf{Question 3.}

\begin{parts}
  \item A body rests in equilibrium upon another fixed body; the portions in contact are spheres of radii $r$ and $R$ respectively, with the line joining centres vertical. If the first body is slightly displaced, find the conditions for stable and unstable equilibrium (the bodies being rough enough to prevent sliding). \pts{6}
  \item A solid sphere rests inside a fixed rough hemispherical bowl of radius twice its own radius. Show that however large a weight is placed at the top of the sphere, the equilibrium is stable. \pts{6}
\end{parts}

\medskip\rule{\linewidth}{0.2pt}

\medskip


\noindent   \textbf{Question 4.}

\begin{parts}
  \item Find the moment of inertia of a circular ring of radius $a$ and mass $M$: (i) about a diameter; (ii) about any tangent in its plane. \pts{6}
  \item If $A$ and $B$ are the moments of inertia of a uniform lamina about two perpendicular axes $OX$ and $OY$ lying in its plane, and $F$ is the product of inertia of the lamina about these lines, show that the principal moments at $O$ are equal to $\dfrac{1}{2}igl(A + B \pm \sqrt{(A-B)^2 + 4F^2}igr)$. \pts{6}
\end{parts}

\medskip\rule{\linewidth}{0.2pt}

\medskip


\noindent   \textbf{Question 5.}

\begin{parts}
  \item Using D'Alembert's principle, derive the equations of motion in Euler's coordinates of a rigid body moving under finite forces in the $OXYZ$ frame. \pts{6}
  \item A uniform vertical circular plate of radius $a$ is capable of revolving about a smooth horizontal axis through its centre. A rough perfectly flexible chain whose mass equals that of the plate and whose length equals its circumference hangs over its rim in equilibrium. If one end is slightly displaced, show that the velocity of the chain when the other end reaches the plate is $\sqrt{\dfrac{nag}{6}}$. \pts{6}
\end{parts}

\medskip\rule{\linewidth}{0.2pt}

\medskip


\noindent   \textbf{Question 6.}

\begin{parts}
  \item A thin uniform rod has one end attached to a smooth hinge and is allowed to fall from a horizontal position. Show that the horizontal reaction on the hinge is greatest when the rod is inclined at $45°$ to the vertical, and that the vertical reaction is then $ frac{11}{8}$ times the weight of the rod. \pts{6}
  \item A cylinder rolls down a smooth inclined plane of inclination $ \theta$, unwrapping as it goes a fine string fixed to the highest point of the plane. Find its acceleration and the tension in the string. \pts{6}
\end{parts}

\medskip\rule{\linewidth}{0.2pt}

\medskip


\noindent   \textbf{Question 7.}

\begin{parts}
  \item If a system moves under finite forces and the geometrical equations of the system do not contain time explicitly, show that the change in the kinetic energy of the system in passing from one configuration to another equals the work done by the forces. \pts{6}
  \item A circular ring of mass $M$ and radius $a$ lies on a smooth horizontal plane and an insect of mass $m$ resting on it starts walking round the ring with uniform velocity $v$ relative to the ring. Show that the centre of the ring describes a circle with angular velocity $\dfrac{mv}{a(M+2m)}$. \pts{6}
\end{parts}

\vfill\rule{\linewidth}{0.4pt}

\begin{center}\small
  \href{https://bhu-exam.vercel.app/}{Click here to visit our website}\\[4pt]
  \href{https://me-aryan.vercel.app/}{Click here to visit our contributor}\\[4pt]
  \href{https://quashan-me.vercel.app/}{Click here to visit our contributor}
\end{center}
\end{document}
